\documentclass[10pt,landscape,letterpaper]{article}
\usepackage[margin=0.38in]{geometry}
\usepackage{booktabs,tabularx,array,xcolor,amssymb,enumitem}
\usepackage{lmodern}
\usepackage[T1]{fontenc}
\usepackage{microtype}
\usepackage{fancyhdr}
\definecolor{navy}{HTML}{173B68}
\definecolor{pale}{HTML}{EAF2F8}
\definecolor{partial}{HTML}{A66A00}
\definecolor{open}{HTML}{657184}
\setlength{\parindent}{0pt}
\setlength{\tabcolsep}{3.2pt}
\renewcommand{\arraystretch}{1.12}
\newcolumntype{Y}{>{\raggedright\arraybackslash}X}
\newcolumntype{P}[1]{>{\raggedright\arraybackslash}p{#1}}
\newcommand{\status}[2]{\textcolor{#1}{#2}}
\pagestyle{fancy}
\fancyhf{}
\fancyfoot[C]{\thepage}
\renewcommand{\headrulewidth}{0pt}
\begin{document}
\color{black}
\begin{flushleft}
{\color{navy}\fontsize{23}{25}\selectfont\bfseries Reviewer Comment Progress Summary\par}
{\large Current evidence reconciliation - 3 September 2026\par}
\end{flushleft}

\vspace{5pt}
\noindent\colorbox{pale}{\parbox{\dimexpr\linewidth-2\fboxsep\relax}{\textbf{Assessment basis.} This update reconciles the rewritten manuscript with the eight-case reactive-film campaign and retained earlier evidence through 3 September 2026. Computed predictions are distinguished from independent validation. Prior numerical scores are retained where applicable; the newly computed film campaign is described without inventing a new grade. Scores are internal assessments, not reviewer-assigned grades.}}

\vspace{7pt}
{\color{navy}\Large Category Assessment\par}
\vspace{2pt}
\small
\begin{tabularx}{\linewidth}{@{}P{1.45in}P{0.84in}P{0.76in}P{0.46in}P{0.56in}Y@{}}
\toprule
\textbf{Category} & \textbf{Comments} & \textbf{Status} & \textbf{Before} & \textbf{Current} & \textbf{Current evidence and remaining work}\\
\midrule
Thermodynamic theory and scope & R1.1-R1.4 & \status{green!50!black}{Addressed} & 1.5, F & 4.8, A & Scope separates reaction equilibrium, transport, and molecular fugacity; the workflow and model context are explicit.\\[2pt]
Parameters, provenance, and sensitivity & R1.5-R1.8; R2.3 & \status{partial}{Partial} & 1.4, F & 3.0, C & The hash-consistent candidate bundle and 131-row regression are retained, but basis admission, independent thermodynamic validation, uncertainty/covariance, complete provenance, and pressure/VLE evidence remain open. The selected parameter candidate is now used in the local film; this use does not establish independent predictive validation.\\[2pt]
Transport-property uncertainty & R1.9 & \status{partial}{Partial} & 2.0, D & 3.5, C+ & Issues 35/42 catalogue source correlations, units, domains, and uncertainty statements. A provisional PSD Onsager closure now passes conservation/current/entropy checks, but species-resolved ionic diffusivities, source-complete cross-friction data, and propagated transport uncertainty remain absent; the eight-case column campaign now reports film fluxes and capture, with estimated mobility.\\[2pt]
Numerical convergence and cost & R1.10-R1.11 & \status{partial}{Partial} & 2.0, D & 3.0, C & All eight hybrid column/film runs meet the declared stopping rule in 5--18 outer updates. Case 5C differs by 0.184 percentage point between two converged relaxation settings. Tighter spatial/outer checks and common-hardware profiling remain incomplete.\\[2pt]
Other amines and transferability & R1.12; R2.6 & \status{partial}{Partial} & 1.0, F & 3.7, B & Solvent-specific requirements are described. No second-solvent demonstration or held-out calibration workflow has been shown.\\[2pt]
Limitations and model context & R1.13-R1.14; R2.2 & \status{green!50!black}{Addressed} & 2.0, D & 4.8, A & Claim boundaries, literature gap, and comparisons with eNRTL, Kent-Eisenberg, CPA, and ePC-SAFT are explicit.\\[2pt]
Formatting and references & R2.1 & \status{green!50!black}{Addressed} & 2.0, D & 5.0, A & Numeric references are in place and citations begin with [1].\\[2pt]
Validation breadth & R2.4 & \status{partial}{Partial} & 2.5, D & 3.4, C+ & Eighteen fixed-chemistry rows pass across K18, K19, and 1C-7C; K20 is rejected. The new film campaign covers eight cases; capture MAE is 6.700 percentage points versus 4.038 for enhancement. Independent cross-facility validation remains absent.\\[2pt]
Operating study and optimization & R2.5 & \status{open}{Open} & 1.0, F & 1.0, F & Still needed: a controlled operating envelope, constraints, objective, sensitivity study, and reproducible optimization.\\[2pt]
Fully predictive reactive ePC-SAFT & Cross-cutting endpoint & \status{open}{Open} & 1.0, F & 1.0, F & A nine-species equilibrium film is now coupled to six-species outer properties. Reactive bulk fugacity and conductance both change. This is computed column transfer, but a fully consistent and independently validated predictive absorber is not established.\\[2pt]
Physical reactive film and column transfer & Issues 30, 36, 42 & \status{partial}{Partial} & 1.0, F & Ungraded & Eight column cases now use the equilibrium-manifold film, with retained flux diagnostics, capture and temperature profiles. Estimated mobility and hybrid chemistry limit interpretation; row-level independent film validation remains open.\\
\bottomrule
\end{tabularx}

\newpage
{\color{navy}\Large Overall Assessment\par}
\vspace{2pt}
\small
\begin{tabularx}{\linewidth}{@{}P{2.45in}P{0.8in}P{1.0in}Y@{}}
\toprule
\textbf{Assessment} & \textbf{Before} & \textbf{Current} & \textbf{Interpretation}\\
\midrule
Overall response to all requested reviewer items & 1.8 / 5, D & 3.1 / 5, C & Theory, scope, limitations, fixed-chemistry reproducibility, source records, and formatting improved. The current evidence does not support predictive reactive chemistry, physical film validation, independent external validation, or optimization.\\[6pt]
Current re-scoped fixed-chemistry benchmark & 1.0 / 5, F & 4.3 / 5, B+/A- & The benchmark has a clear claim boundary and stronger reproducibility evidence: 18 accepted rows across K18, K19, and 1C-7C for both closures, with K20 retained as rejected.\\[6pt]
Fully predictive reactive ePC-SAFT endpoint & 1.0 / 5, F & 1.0 / 5, F & The nine-species local film now supplies column transfer. Estimated mobility, six-species outer properties, the R5 temperature extrapolation, and absent independent validation prevent a fully predictive absorber claim.\\[6pt]
Physical reactive-film / column calculation & 1.0 / 5, F & Ungraded & The current campaign establishes computed film fluxes and eight converged column predictions. It does not isolate transfer at fixed reactive chemistry and does not demonstrate superior capture agreement or independently validated transport.\\
\bottomrule
\end{tabularx}

\vspace{9pt}
{\color{navy}\Large Verified Updates and Claim Boundaries\par}
\vspace{3pt}
\begin{itemize}[leftmargin=0.18in,itemsep=3pt,topsep=2pt]
\item \textbf{Manuscript and presentation:} the approved ePC-SAFT-centered article now explains the nine-species local equilibrium, projected mobility, gas-film interface balance, six-species outer coupling, and separate column BVP. Tables, vector diagrams, and grayscale chart distinctions are updated.
\item \textbf{Completed column calculations:} K18, K19 and 1C--6C all meet the declared outer stopping rule. Film capture MAE is 6.700 percentage points, compared with 4.038 for enhancement; the film does not improve aggregate capture agreement.
\item \textbf{Temperature evidence:} six C-case liquid-temperature tap RMSE values range from 3.73 to 8.05 K, with a mean of 5.65 K. Profiles and exact metrics are retained alongside capture results.
\item \textbf{Numerical limits:} the maximum collocation RMS residual is 0.0661 at tolerance 0.1; a small stationary-flux residual is not a mesh-independence result. Two converged 5C relaxation settings differ by 0.184 percentage point.
\item \textbf{Physical interpretation:} both reactive bulk fugacity and conductance change between enhancement and film, so the current hybrid comparison cannot isolate transfer alone. Species diffusivities are estimates. Of 500 retained local evaluations, 369 exceed the 323.15 K source ceiling for reaction R5.
\item \textbf{Remaining comparisons:} a common-chemistry enhancement/film study, a Henry/eNRTL/ePC-SAFT comparison, and a solver comparison using the same reactive-film equations remain open. Earlier solver results apply to the enhancement model. Historical blocked source gates remain historical evidence and are not reclassified by successful numerical calculations.
\end{itemize}

\vspace{4pt}
\noindent\fcolorbox{navy}{pale}{\parbox{\dimexpr\linewidth-2\fboxsep-2\fboxrule\relax}{\textbf{Current conclusion.} The film is now implemented and exercised in eight column cases, and its limitations are explicit in the manuscript. This is material computational progress. Independent predictive validation, consistent reactive chemistry throughout the column, and the missing controlled comparisons remain unfinished.}}

\vfill
\scriptsize\textit{Sources: current reactive-film final capture, temperature, node and relaxation tables; immutable Engine commit 8007a708; earlier fixed-chemistry benchmark and source reports. Scientific manuscript revision dated 3 September 2026.}
\end{document}
